
\documentclass{article}
\usepackage{colortbl}
\usepackage{makecell}
\usepackage{multirow}
\usepackage{supertabular}

\begin{document}

\newcounter{utterance}

\centering \large Interaction Transcript for game `clean\_up', experiment `0\_easy\_5obj\_sl', episode 1 with nemotron{-}3{-}ultra.
\vspace{24pt}

{ \footnotesize  \setcounter{utterance}{1}
\setlength{\tabcolsep}{0pt}
\begin{supertabular}{c@{$\;$}|p{.15\linewidth}@{}p{.15\linewidth}p{.15\linewidth}p{.15\linewidth}p{.15\linewidth}p{.15\linewidth}}
    \# & $\;$A & \multicolumn{4}{c}{Game Master} & $\;\:$B\\
    \hline

    \theutterance \stepcounter{utterance}  
    & & \multicolumn{4}{p{0.6\linewidth}}{
        \cellcolor[rgb]{0.9,0.9,0.9}{
            \makecell[{{p{\linewidth}}}]{
                \texttt{\tiny{[P1$\langle$GM]}}
                \texttt{Jaz sem vaš vodja igre, vi pa igrate skupinsko igro z naslednjo mrežo kot igralno ploščo:} \\
\\ 
\texttt{```} \\
\texttt{ 1234567} \\
\texttt{╔═══════╗} \\
\texttt{║◌◌◌◌◌◌◌║ 1} \\
\texttt{║◌◌WH◌◌◌║ 2} \\
\texttt{║◌C◌◌◌┌─╢ 3} \\
\texttt{║◌◌◌┌─┤◌║ 4} \\
\texttt{║◌◌I│◌├─╢ 5} \\
\texttt{╟───┼─┘◌║ 6} \\
\texttt{║◌◌T│◌◌◌║ 7} \\
\texttt{╚═══╧═══╝} \\
\\ 
\texttt{```} \\
\\ 
\texttt{* Na zgornjem robu so prikazane koordinate x, ki naraščajo proti desni, na desnem robu pa koordinate y, ki naraščajo navzdol.} \\
\texttt{* Na mrežo so naključno postavljeni naslednji predmeti: 'W', 'I', 'T', 'C', 'H'.} \\
\\ 
\\ 
\texttt{Drugi igralec vidi različico igralne plošče, na kateri so predmeti postavljeni na različna naključna mesta. Vi ne morete videti plošče drugega igralca, on pa ne more videti vaše.} \\
\\ 
\\ 
\texttt{**Cilj:**} \\
\\ 
\texttt{Oba igralca morata premikati predmete na svojem ozadju tako, da se enaki predmeti znajdejo na istih koordinatah. Z drugim igralcem se morate dogovoriti o skupni ciljni konfiguraciji.} \\
\\ 
\\ 
\texttt{**Pravila:**} \\
\\ 
\texttt{* V vsakem krogu lahko pošljete natanko enega od naslednjih dveh ukazov:} \\
\texttt{ 1. `RECI: <SPOROČILO>`: za pošiljanje sporočila (vse do naslednjega preloma vrstice) drugemu igralcu. Posredoval ga bom vašemu partnerju.} \\
\texttt{2. `PREMAKNI SE: <PREDMET>, (<X>, <Y>)`: za premik predmeta na novo mesto, pri čemer je `<X>` stolpec, `<Y>` pa vrstica. Obvestil vas bom, ali je bila vaša premestitev uspešna ali ne.} \\
\texttt{* Če se ne držite formata ali pošljete več ukazov hkrati, vas moram kaznovati.} \\
\texttt{* Če oba igralca zbereta več kot 12 kazni, oba izgubita igro.} \\
\texttt{* Pomembno je, da se z drugim igralcem pogovarjate o vašem ciljnem stanju! Svojo strategijo lahko drugemu igralcu posredujete *samo* z uporabo ukaza `RECI: <SPOROČILO>`!} \\
\\ 
\\ 
\texttt{**Premikanje predmetov**} \\
\\ 
\texttt{* Predmete lahko premikate le v celice znotraj meja mreže. Ciljna celica mora biti prazna, tj. vsebovati mora samo simbol '◌'.} \\
\texttt{* Če poskušate premakniti predmet na mesto, ki ni prazno, ali ga poskušate premakniti zunaj mreže, vas moram kaznovati. Dobili boste še en poskus.} \\
\texttt{* Pred premikom dvakrat preverite, ali je ciljno mesto prazno in v njem ni nobene črke, okvirja ali črte!} \\
\\ 
\\ 
\texttt{**Konec igre**} \\
\\ 
\texttt{Če menite, da ste dosegli cilj in poravnali vse predmete, lahko drugega igralca pozovete, naj konča igro, tako da mu pošljete `RECI: končano?`. Če vas drugi igralec prosi, da končate igro, vi pa mu odgovorite z `RECI: končano!`, se igra konča.} \\
\\ 
\texttt{Oba igralca zmagata, če se igra konča v 20 krogih, pri čemer je en krog definiran tako, da dva igralca pošljeta vsak po en veljaven ukaz.} \\
\\ 
\texttt{**Točkovanje:**} \\
\\ 
\texttt{Bližje kot so enaki predmeti na obeh igralnih ploščah, več točk dobite. Kazni zmanjšajo število točk. Ali lahko premagate rekord?} \\
\\ 
\\ 
\texttt{Pošljite sporočilo drugemu igralcu, da začnete igro!} \\
            }
        }
    }
    & \\ \\

    \theutterance \stepcounter{utterance}  
    & \multicolumn{4}{p{0.6\linewidth}}{
        \cellcolor[rgb]{0.9,0.9,0.9}{
            \makecell[{{p{\linewidth}}}]{
                \texttt{\tiny{[P1$\rangle$GM]}}
                \texttt{RECI: Pozdravljen! Poglejmo si začetne pozicije predmetov na moji plošči: W (2,2), H (4,2), C (2,3), I (3,5), T (3,7). Kje so tvoji predmeti postavljeni?} \\
            }
        }
    }
    & & \\ \\

    \theutterance \stepcounter{utterance}  
    & & & \multicolumn{2}{p{0.3\linewidth}}{
        \cellcolor[rgb]{0.9,0.9,0.9}{
            \makecell[{{p{\linewidth}}}]{
                \texttt{\tiny{[GM$|$GM]}}
                \texttt{RECI: Pozdravljen! Poglejmo si začetne pozicije predmetov na moji plošči: W (2,2), H (4,2), C (2,3), I (3,5), T (3,7). Kje so tvoji predmeti postavljeni?} \\
            }
        }
    }
    & & \\ \\

    \theutterance \stepcounter{utterance}  
    & & \multicolumn{4}{p{0.6\linewidth}}{
        \cellcolor[rgb]{0.9,0.9,0.9}{
            \makecell[{{p{\linewidth}}}]{
                \texttt{\tiny{[GM$\rangle$P2]}}
                \texttt{Jaz sem vaš vodja igre, vi pa igrate skupinsko igro z naslednjo mrežo kot igralno ploščo:} \\
\\ 
\texttt{```} \\
\texttt{ 1234567} \\
\texttt{╔═══════╗} \\
\texttt{║◌◌◌◌◌◌◌║ 1} \\
\texttt{║◌◌H◌◌◌◌║ 2} \\
\texttt{║◌WC◌◌┌─╢ 3} \\
\texttt{║◌◌◌┌─┤◌║ 4} \\
\texttt{║T◌◌│I├─╢ 5} \\
\texttt{╟───┼─┘◌║ 6} \\
\texttt{║◌◌◌│◌◌◌║ 7} \\
\texttt{╚═══╧═══╝} \\
\\ 
\texttt{```} \\
\\ 
\texttt{* Na zgornjem robu so prikazane koordinate x, ki naraščajo proti desni, na desnem robu pa koordinate y, ki naraščajo navzdol.} \\
\texttt{* Na mrežo so naključno postavljeni naslednji predmeti: 'W', 'I', 'T', 'C', 'H'.} \\
\\ 
\\ 
\texttt{Drugi igralec vidi različico igralne plošče, na kateri so predmeti postavljeni na različna naključna mesta. Vi ne morete videti plošče drugega igralca, on pa ne more videti vaše.} \\
\\ 
\\ 
\texttt{**Cilj:**} \\
\\ 
\texttt{Oba igralca morata premikati predmete na svojem ozadju tako, da se enaki predmeti znajdejo na istih koordinatah. Z drugim igralcem se morate dogovoriti o skupni ciljni konfiguraciji.} \\
\\ 
\\ 
\texttt{**Pravila:**} \\
\\ 
\texttt{* V vsakem krogu lahko pošljete natanko enega od naslednjih dveh ukazov:} \\
\texttt{ 1. `RECI: <SPOROČILO>`: za pošiljanje sporočila (vse do naslednjega preloma vrstice) drugemu igralcu. Posredoval ga bom vašemu partnerju.} \\
\texttt{2. `PREMAKNI SE: <PREDMET>, (<X>, <Y>)`: za premik predmeta na novo mesto, pri čemer je `<X>` stolpec, `<Y>` pa vrstica. Obvestil vas bom, ali je bila vaša premestitev uspešna ali ne.} \\
\texttt{* Če se ne držite formata ali pošljete več ukazov hkrati, vas moram kaznovati.} \\
\texttt{* Če oba igralca zbereta več kot 12 kazni, oba izgubita igro.} \\
\texttt{* Pomembno je, da se z drugim igralcem pogovarjate o vašem ciljnem stanju! Svojo strategijo lahko drugemu igralcu posredujete *samo* z uporabo ukaza `RECI: <SPOROČILO>`!} \\
\\ 
\\ 
\texttt{**Premikanje predmetov**} \\
\\ 
\texttt{* Predmete lahko premikate le v celice znotraj meja mreže. Ciljna celica mora biti prazna, tj. vsebovati mora samo simbol '◌'.} \\
\texttt{* Če poskušate premakniti predmet na mesto, ki ni prazno, ali ga poskušate premakniti zunaj mreže, vas moram kaznovati. Dobili boste še en poskus.} \\
\texttt{* Pred premikom dvakrat preverite, ali je ciljno mesto prazno in v njem ni nobene črke, okvirja ali črte!} \\
\\ 
\\ 
\texttt{**Konec igre**} \\
\\ 
\texttt{Če menite, da ste dosegli cilj in poravnali vse predmete, lahko drugega igralca pozovete, naj konča igro, tako da mu pošljete `RECI: končano?`. Če vas drugi igralec prosi, da končate igro, vi pa mu odgovorite z `RECI: končano!`, se igra konča.} \\
\\ 
\texttt{Oba igralca zmagata, če se igra konča v 20 krogih, pri čemer je en krog definiran tako, da dva igralca pošljeta vsak po en veljaven ukaz.} \\
\\ 
\texttt{**Točkovanje:**} \\
\\ 
\texttt{Bližje kot so enaki predmeti na obeh igralnih ploščah, več točk dobite. Kazni zmanjšajo število točk. Ali lahko premagate rekord?} \\
\\ 
\texttt{Drugi igralec je začel igro s tem sporočilom:} \\
\\ 
\texttt{"Pozdravljen! Poglejmo si začetne pozicije predmetov na moji plošči: W (2,2), H (4,2), C (2,3), I (3,5), T (3,7). Kje so tvoji predmeti postavljeni?"} \\
\\ 
\texttt{Kakšen je vaš prvi ukaz?} \\
            }
        }
    }
    & \\ \\

    \theutterance \stepcounter{utterance}  
    & & & \multicolumn{4}{p{0.6\linewidth}}{
        \cellcolor[rgb]{0.9,0.9,0.9}{
            \makecell[{{p{\linewidth}}}]{
                \texttt{\tiny{[GM$\langle$P2]}}
                \texttt{RECI: Na moji plošči so predmeti na: W (3,3), H (3,2), C (3,4), I (4,5), T (1,5). Predlagam, da jih premaknemo na sledeče koordinate: W (1,1), H (2,2), C (3,3), I (4,4), T (5,5). Ali strinjate?} \\
            }
        }
    }
     \\ \\

    \theutterance \stepcounter{utterance}  
    & & & \multicolumn{2}{p{0.3\linewidth}}{
        \cellcolor[rgb]{0.9,0.9,0.9}{
            \makecell[{{p{\linewidth}}}]{
                \texttt{\tiny{[GM$|$GM]}}
                \texttt{RECI: Na moji plošči so predmeti na: W (3,3), H (3,2), C (3,4), I (4,5), T (1,5). Predlagam, da jih premaknemo na sledeče koordinate: W (1,1), H (2,2), C (3,3), I (4,4), T (5,5). Ali strinjate?} \\
            }
        }
    }
    & & \\ \\

    \theutterance \stepcounter{utterance}  
    & & \multicolumn{4}{p{0.6\linewidth}}{
        \cellcolor[rgb]{0.9,0.9,0.9}{
            \makecell[{{p{\linewidth}}}]{
                \texttt{\tiny{[P1$\langle$GM]}}
                \texttt{Vaše sporočilo je bilo posredovano drugemu igralcu.} \\
\\ 
\\ 
\texttt{Trenutno igrate krog 1 od največ 20.} \\
\texttt{Skupaj ste zbrali 0 od 12 kazni.} \\
\texttt{Drugi igralec je poslal to sporočilo:} \\
\\ 
\\ 
\texttt{"Na moji plošči so predmeti na: W (3,3), H (3,2), C (3,4), I (4,5), T (1,5). Predlagam, da jih premaknemo na sledeče koordinate: W (1,1), H (2,2), C (3,3), I (4,4), T (5,5). Ali strinjate?"} \\
\\ 
\\ 
\texttt{Kakšen je vaš naslednji ukaz?} \\
            }
        }
    }
    & \\ \\

    \theutterance \stepcounter{utterance}  
    & \multicolumn{4}{p{0.6\linewidth}}{
        \cellcolor[rgb]{0.9,0.9,0.9}{
            \makecell[{{p{\linewidth}}}]{
                \texttt{\tiny{[P1$\rangle$GM]}}
                \texttt{RECI: Strinjam se s predlagano ciljno konfiguracijo: W (1,1), H (2,2), C (3,3), I (4,4), T (5,5). Začenjam s premiki.} \\
            }
        }
    }
    & & \\ \\

    \theutterance \stepcounter{utterance}  
    & & & \multicolumn{2}{p{0.3\linewidth}}{
        \cellcolor[rgb]{0.9,0.9,0.9}{
            \makecell[{{p{\linewidth}}}]{
                \texttt{\tiny{[GM$|$GM]}}
                \texttt{RECI: Strinjam se s predlagano ciljno konfiguracijo: W (1,1), H (2,2), C (3,3), I (4,4), T (5,5). Začenjam s premiki.} \\
            }
        }
    }
    & & \\ \\

    \theutterance \stepcounter{utterance}  
    & & \multicolumn{4}{p{0.6\linewidth}}{
        \cellcolor[rgb]{0.9,0.9,0.9}{
            \makecell[{{p{\linewidth}}}]{
                \texttt{\tiny{[GM$\rangle$P2]}}
                \texttt{Vaše sporočilo je bilo posredovano drugemu igralcu.} \\
\\ 
\\ 
\texttt{Trenutno igrate krog 2 od največ 20.} \\
\texttt{Skupaj ste zbrali 0 od 12 kazni.} \\
\texttt{Drugi igralec je poslal to sporočilo:} \\
\\ 
\\ 
\texttt{"Strinjam se s predlagano ciljno konfiguracijo: W (1,1), H (2,2), C (3,3), I (4,4), T (5,5). Začenjam s premiki."} \\
\\ 
\\ 
\texttt{Kakšen je vaš naslednji ukaz?} \\
            }
        }
    }
    & \\ \\

    \theutterance \stepcounter{utterance}  
    & & & \multicolumn{4}{p{0.6\linewidth}}{
        \cellcolor[rgb]{0.9,0.9,0.9}{
            \makecell[{{p{\linewidth}}}]{
                \texttt{\tiny{[GM$\langle$P2]}}
                \texttt{PREMAKNI SE: W, (1,1)} \\
            }
        }
    }
     \\ \\

    \theutterance \stepcounter{utterance}  
    & & & \multicolumn{2}{p{0.3\linewidth}}{
        \cellcolor[rgb]{0.9,0.9,0.9}{
            \makecell[{{p{\linewidth}}}]{
                \texttt{\tiny{[GM$|$GM]}}
                \texttt{PREMAKNI SE: W, (1,1)} \\
            }
        }
    }
    & & \\ \\

    \theutterance \stepcounter{utterance}  
    & & & \multicolumn{2}{p{0.3\linewidth}}{
        \cellcolor[rgb]{0.9,0.9,0.9}{
            \makecell[{{p{\linewidth}}}]{
                \texttt{\tiny{[GM$|$GM]}}
                \texttt{W moved to (1, 1)} \\
            }
        }
    }
    & & \\ \\

    \theutterance \stepcounter{utterance}  
    & & \multicolumn{4}{p{0.6\linewidth}}{
        \cellcolor[rgb]{0.9,0.9,0.9}{
            \makecell[{{p{\linewidth}}}]{
                \texttt{\tiny{[P1$\langle$GM]}}
                \texttt{Vaše sporočilo je bilo posredovano drugemu igralcu.} \\
\\ 
\\ 
\texttt{Trenutno igrate krog 2 od največ 20.} \\
\texttt{Skupaj ste zbrali 0 od 12 kazni.} \\
\texttt{Drugi igralec je premaknil predmet na svoji mreži.} \\
\\ 
\\ 
\texttt{Kakšen je vaš naslednji ukaz?} \\
            }
        }
    }
    & \\ \\

    \theutterance \stepcounter{utterance}  
    & \multicolumn{4}{p{0.6\linewidth}}{
        \cellcolor[rgb]{0.9,0.9,0.9}{
            \makecell[{{p{\linewidth}}}]{
                \texttt{\tiny{[P1$\rangle$GM]}}
                \texttt{PREMAKNI SE: W, (1, 1)} \\
            }
        }
    }
    & & \\ \\

    \theutterance \stepcounter{utterance}  
    & & & \multicolumn{2}{p{0.3\linewidth}}{
        \cellcolor[rgb]{0.9,0.9,0.9}{
            \makecell[{{p{\linewidth}}}]{
                \texttt{\tiny{[GM$|$GM]}}
                \texttt{PREMAKNI SE: W, (1, 1)} \\
            }
        }
    }
    & & \\ \\

    \theutterance \stepcounter{utterance}  
    & & & \multicolumn{2}{p{0.3\linewidth}}{
        \cellcolor[rgb]{0.9,0.9,0.9}{
            \makecell[{{p{\linewidth}}}]{
                \texttt{\tiny{[GM$|$GM]}}
                \texttt{W moved to (1, 1)} \\
            }
        }
    }
    & & \\ \\

    \theutterance \stepcounter{utterance}  
    & & \multicolumn{4}{p{0.6\linewidth}}{
        \cellcolor[rgb]{0.9,0.9,0.9}{
            \makecell[{{p{\linewidth}}}]{
                \texttt{\tiny{[GM$\rangle$P2]}}
                \texttt{Uspešno premaknjen 'W' v (1,1). Vaša posodobljena mreža je videti takole:} \\
\texttt{```} \\
\texttt{ 1234567} \\
\texttt{╔═══════╗} \\
\texttt{║W◌◌◌◌◌◌║ 1} \\
\texttt{║◌◌H◌◌◌◌║ 2} \\
\texttt{║◌◌C◌◌┌─╢ 3} \\
\texttt{║◌◌◌┌─┤◌║ 4} \\
\texttt{║T◌◌│I├─╢ 5} \\
\texttt{╟───┼─┘◌║ 6} \\
\texttt{║◌◌◌│◌◌◌║ 7} \\
\texttt{╚═══╧═══╝} \\
\texttt{```} \\
\texttt{Trenutno igrate krog 3 od največ 20.} \\
\texttt{Skupaj ste zbrali 0 od 12 kazni.} \\
\texttt{Drugi igralec je premaknil predmet na svoji mreži.} \\
\\ 
\\ 
\texttt{Kakšen je vaš naslednji ukaz?} \\
            }
        }
    }
    & \\ \\

    \theutterance \stepcounter{utterance}  
    & & & \multicolumn{4}{p{0.6\linewidth}}{
        \cellcolor[rgb]{0.9,0.9,0.9}{
            \makecell[{{p{\linewidth}}}]{
                \texttt{\tiny{[GM$\langle$P2]}}
                \texttt{PREMAKNI SE: H, (2,2)} \\
            }
        }
    }
     \\ \\

    \theutterance \stepcounter{utterance}  
    & & & \multicolumn{2}{p{0.3\linewidth}}{
        \cellcolor[rgb]{0.9,0.9,0.9}{
            \makecell[{{p{\linewidth}}}]{
                \texttt{\tiny{[GM$|$GM]}}
                \texttt{PREMAKNI SE: H, (2,2)} \\
            }
        }
    }
    & & \\ \\

    \theutterance \stepcounter{utterance}  
    & & & \multicolumn{2}{p{0.3\linewidth}}{
        \cellcolor[rgb]{0.9,0.9,0.9}{
            \makecell[{{p{\linewidth}}}]{
                \texttt{\tiny{[GM$|$GM]}}
                \texttt{H moved to (2, 2)} \\
            }
        }
    }
    & & \\ \\

    \theutterance \stepcounter{utterance}  
    & & \multicolumn{4}{p{0.6\linewidth}}{
        \cellcolor[rgb]{0.9,0.9,0.9}{
            \makecell[{{p{\linewidth}}}]{
                \texttt{\tiny{[P1$\langle$GM]}}
                \texttt{Uspešno premaknjen 'W' v (1,1). Vaša posodobljena mreža je videti takole:} \\
\texttt{```} \\
\texttt{ 1234567} \\
\texttt{╔═══════╗} \\
\texttt{║W◌◌◌◌◌◌║ 1} \\
\texttt{║◌◌◌H◌◌◌║ 2} \\
\texttt{║◌C◌◌◌┌─╢ 3} \\
\texttt{║◌◌◌┌─┤◌║ 4} \\
\texttt{║◌◌I│◌├─╢ 5} \\
\texttt{╟───┼─┘◌║ 6} \\
\texttt{║◌◌T│◌◌◌║ 7} \\
\texttt{╚═══╧═══╝} \\
\texttt{```} \\
\texttt{Trenutno igrate krog 3 od največ 20.} \\
\texttt{Skupaj ste zbrali 0 od 12 kazni.} \\
\texttt{Drugi igralec je premaknil predmet na svoji mreži.} \\
\\ 
\\ 
\texttt{Kakšen je vaš naslednji ukaz?} \\
            }
        }
    }
    & \\ \\

    \theutterance \stepcounter{utterance}  
    & \multicolumn{4}{p{0.6\linewidth}}{
        \cellcolor[rgb]{0.9,0.9,0.9}{
            \makecell[{{p{\linewidth}}}]{
                \texttt{\tiny{[P1$\rangle$GM]}}
                \texttt{PREMAKNI SE: H, (2, 2)} \\
            }
        }
    }
    & & \\ \\

    \theutterance \stepcounter{utterance}  
    & & & \multicolumn{2}{p{0.3\linewidth}}{
        \cellcolor[rgb]{0.9,0.9,0.9}{
            \makecell[{{p{\linewidth}}}]{
                \texttt{\tiny{[GM$|$GM]}}
                \texttt{PREMAKNI SE: H, (2, 2)} \\
            }
        }
    }
    & & \\ \\

    \theutterance \stepcounter{utterance}  
    & & & \multicolumn{2}{p{0.3\linewidth}}{
        \cellcolor[rgb]{0.9,0.9,0.9}{
            \makecell[{{p{\linewidth}}}]{
                \texttt{\tiny{[GM$|$GM]}}
                \texttt{H moved to (2, 2)} \\
            }
        }
    }
    & & \\ \\

    \theutterance \stepcounter{utterance}  
    & & \multicolumn{4}{p{0.6\linewidth}}{
        \cellcolor[rgb]{0.9,0.9,0.9}{
            \makecell[{{p{\linewidth}}}]{
                \texttt{\tiny{[GM$\rangle$P2]}}
                \texttt{Uspešno premaknjen 'H' v (2,2). Vaša posodobljena mreža je videti takole:} \\
\texttt{```} \\
\texttt{ 1234567} \\
\texttt{╔═══════╗} \\
\texttt{║W◌◌◌◌◌◌║ 1} \\
\texttt{║◌H◌◌◌◌◌║ 2} \\
\texttt{║◌◌C◌◌┌─╢ 3} \\
\texttt{║◌◌◌┌─┤◌║ 4} \\
\texttt{║T◌◌│I├─╢ 5} \\
\texttt{╟───┼─┘◌║ 6} \\
\texttt{║◌◌◌│◌◌◌║ 7} \\
\texttt{╚═══╧═══╝} \\
\texttt{```} \\
\texttt{Trenutno igrate krog 4 od največ 20.} \\
\texttt{Skupaj ste zbrali 0 od 12 kazni.} \\
\texttt{Drugi igralec je premaknil predmet na svoji mreži.} \\
\\ 
\\ 
\texttt{Kakšen je vaš naslednji ukaz?} \\
            }
        }
    }
    & \\ \\

    \theutterance \stepcounter{utterance}  
    & & & \multicolumn{4}{p{0.6\linewidth}}{
        \cellcolor[rgb]{0.9,0.9,0.9}{
            \makecell[{{p{\linewidth}}}]{
                \texttt{\tiny{[GM$\langle$P2]}}
                \texttt{PREMAKNI SE: C, (3,3)} \\
            }
        }
    }
     \\ \\

    \theutterance \stepcounter{utterance}  
    & & & \multicolumn{2}{p{0.3\linewidth}}{
        \cellcolor[rgb]{0.9,0.9,0.9}{
            \makecell[{{p{\linewidth}}}]{
                \texttt{\tiny{[GM$|$GM]}}
                \texttt{PREMAKNI SE: C, (3,3)} \\
            }
        }
    }
    & & \\ \\

    \theutterance \stepcounter{utterance}  
    & & & \multicolumn{2}{p{0.3\linewidth}}{
        \cellcolor[rgb]{0.9,0.9,0.9}{
            \makecell[{{p{\linewidth}}}]{
                \texttt{\tiny{[GM$|$GM]}}
                \texttt{Kazen: (3,3) ni prazen, vendar vsebuje 'C'.} \\
\\ 
\texttt{Skupaj ste zbrali 1 od 12 kazni.} \\
\texttt{Prosim, poskusite znova!} \\
            }
        }
    }
    & & \\ \\

    \theutterance \stepcounter{utterance}  
    & & \multicolumn{4}{p{0.6\linewidth}}{
        \cellcolor[rgb]{0.9,0.9,0.9}{
            \makecell[{{p{\linewidth}}}]{
                \texttt{\tiny{[GM$\rangle$P2]}}
                \texttt{Kazen: (3,3) ni prazen, vendar vsebuje 'C'.} \\
\\ 
\texttt{Skupaj ste zbrali 1 od 12 kazni.} \\
\texttt{Prosim, poskusite znova!} \\
            }
        }
    }
    & \\ \\

    \theutterance \stepcounter{utterance}  
    & & & \multicolumn{4}{p{0.6\linewidth}}{
        \cellcolor[rgb]{0.9,0.9,0.9}{
            \makecell[{{p{\linewidth}}}]{
                \texttt{\tiny{[GM$\langle$P2]}}
                \texttt{RECI: C je že na (3,3), zato ga ne premikam. Nadaljujem s T.} \\
            }
        }
    }
     \\ \\

    \theutterance \stepcounter{utterance}  
    & & & \multicolumn{2}{p{0.3\linewidth}}{
        \cellcolor[rgb]{0.9,0.9,0.9}{
            \makecell[{{p{\linewidth}}}]{
                \texttt{\tiny{[GM$|$GM]}}
                \texttt{RECI: C je že na (3,3), zato ga ne premikam. Nadaljujem s T.} \\
            }
        }
    }
    & & \\ \\

    \theutterance \stepcounter{utterance}  
    & & \multicolumn{4}{p{0.6\linewidth}}{
        \cellcolor[rgb]{0.9,0.9,0.9}{
            \makecell[{{p{\linewidth}}}]{
                \texttt{\tiny{[P1$\langle$GM]}}
                \texttt{Uspešno premaknjen 'H' v (2,2). Vaša posodobljena mreža je videti takole:} \\
\texttt{```} \\
\texttt{ 1234567} \\
\texttt{╔═══════╗} \\
\texttt{║W◌◌◌◌◌◌║ 1} \\
\texttt{║◌H◌◌◌◌◌║ 2} \\
\texttt{║◌C◌◌◌┌─╢ 3} \\
\texttt{║◌◌◌┌─┤◌║ 4} \\
\texttt{║◌◌I│◌├─╢ 5} \\
\texttt{╟───┼─┘◌║ 6} \\
\texttt{║◌◌T│◌◌◌║ 7} \\
\texttt{╚═══╧═══╝} \\
\texttt{```} \\
\texttt{Trenutno igrate krog 4 od največ 20.} \\
\texttt{Skupaj ste zbrali 1 od 12 kazni.} \\
\texttt{Drugi igralec je poslal to sporočilo:} \\
\\ 
\\ 
\texttt{"C je že na (3,3), zato ga ne premikam. Nadaljujem s T."} \\
\\ 
\\ 
\texttt{Kakšen je vaš naslednji ukaz?} \\
            }
        }
    }
    & \\ \\

    \theutterance \stepcounter{utterance}  
    & \multicolumn{4}{p{0.6\linewidth}}{
        \cellcolor[rgb]{0.9,0.9,0.9}{
            \makecell[{{p{\linewidth}}}]{
                \texttt{\tiny{[P1$\rangle$GM]}}
                \texttt{PREMAKNI SE: C, (3, 3)} \\
            }
        }
    }
    & & \\ \\

    \theutterance \stepcounter{utterance}  
    & & & \multicolumn{2}{p{0.3\linewidth}}{
        \cellcolor[rgb]{0.9,0.9,0.9}{
            \makecell[{{p{\linewidth}}}]{
                \texttt{\tiny{[GM$|$GM]}}
                \texttt{PREMAKNI SE: C, (3, 3)} \\
            }
        }
    }
    & & \\ \\

    \theutterance \stepcounter{utterance}  
    & & & \multicolumn{2}{p{0.3\linewidth}}{
        \cellcolor[rgb]{0.9,0.9,0.9}{
            \makecell[{{p{\linewidth}}}]{
                \texttt{\tiny{[GM$|$GM]}}
                \texttt{C moved to (3, 3)} \\
            }
        }
    }
    & & \\ \\

    \theutterance \stepcounter{utterance}  
    & & \multicolumn{4}{p{0.6\linewidth}}{
        \cellcolor[rgb]{0.9,0.9,0.9}{
            \makecell[{{p{\linewidth}}}]{
                \texttt{\tiny{[GM$\rangle$P2]}}
                \texttt{Vaše sporočilo je bilo posredovano drugemu igralcu.} \\
\\ 
\\ 
\texttt{Trenutno igrate krog 5 od največ 20.} \\
\texttt{Skupaj ste zbrali 1 od 12 kazni.} \\
\texttt{Drugi igralec je premaknil predmet na svoji mreži.} \\
\\ 
\\ 
\texttt{Kakšen je vaš naslednji ukaz?} \\
            }
        }
    }
    & \\ \\

    \theutterance \stepcounter{utterance}  
    & & & \multicolumn{4}{p{0.6\linewidth}}{
        \cellcolor[rgb]{0.9,0.9,0.9}{
            \makecell[{{p{\linewidth}}}]{
                \texttt{\tiny{[GM$\langle$P2]}}
                \texttt{PREMAKNI SE: T, (5,5)} \\
            }
        }
    }
     \\ \\

    \theutterance \stepcounter{utterance}  
    & & & \multicolumn{2}{p{0.3\linewidth}}{
        \cellcolor[rgb]{0.9,0.9,0.9}{
            \makecell[{{p{\linewidth}}}]{
                \texttt{\tiny{[GM$|$GM]}}
                \texttt{PREMAKNI SE: T, (5,5)} \\
            }
        }
    }
    & & \\ \\

    \theutterance \stepcounter{utterance}  
    & & & \multicolumn{2}{p{0.3\linewidth}}{
        \cellcolor[rgb]{0.9,0.9,0.9}{
            \makecell[{{p{\linewidth}}}]{
                \texttt{\tiny{[GM$|$GM]}}
                \texttt{Kazen: (5,5) ni prazen, vendar vsebuje 'I'.} \\
\\ 
\texttt{Skupaj ste zbrali 2 od 12 kazni.} \\
\texttt{Prosim, poskusite znova!} \\
            }
        }
    }
    & & \\ \\

    \theutterance \stepcounter{utterance}  
    & & \multicolumn{4}{p{0.6\linewidth}}{
        \cellcolor[rgb]{0.9,0.9,0.9}{
            \makecell[{{p{\linewidth}}}]{
                \texttt{\tiny{[GM$\rangle$P2]}}
                \texttt{Kazen: (5,5) ni prazen, vendar vsebuje 'I'.} \\
\\ 
\texttt{Skupaj ste zbrali 2 od 12 kazni.} \\
\texttt{Prosim, poskusite znova!} \\
            }
        }
    }
    & \\ \\

    \theutterance \stepcounter{utterance}  
    & & & \multicolumn{4}{p{0.6\linewidth}}{
        \cellcolor[rgb]{0.9,0.9,0.9}{
            \makecell[{{p{\linewidth}}}]{
                \texttt{\tiny{[GM$\langle$P2]}}
                \texttt{PREMAKNI SE: T, (5,3)} \\
            }
        }
    }
     \\ \\

    \theutterance \stepcounter{utterance}  
    & & & \multicolumn{2}{p{0.3\linewidth}}{
        \cellcolor[rgb]{0.9,0.9,0.9}{
            \makecell[{{p{\linewidth}}}]{
                \texttt{\tiny{[GM$|$GM]}}
                \texttt{PREMAKNI SE: T, (5,3)} \\
            }
        }
    }
    & & \\ \\

    \theutterance \stepcounter{utterance}  
    & & & \multicolumn{2}{p{0.3\linewidth}}{
        \cellcolor[rgb]{0.9,0.9,0.9}{
            \makecell[{{p{\linewidth}}}]{
                \texttt{\tiny{[GM$|$GM]}}
                \texttt{T moved to (5, 3)} \\
            }
        }
    }
    & & \\ \\

    \theutterance \stepcounter{utterance}  
    & & \multicolumn{4}{p{0.6\linewidth}}{
        \cellcolor[rgb]{0.9,0.9,0.9}{
            \makecell[{{p{\linewidth}}}]{
                \texttt{\tiny{[P1$\langle$GM]}}
                \texttt{Uspešno premaknjen 'C' v (3,3). Vaša posodobljena mreža je videti takole:} \\
\texttt{```} \\
\texttt{ 1234567} \\
\texttt{╔═══════╗} \\
\texttt{║W◌◌◌◌◌◌║ 1} \\
\texttt{║◌H◌◌◌◌◌║ 2} \\
\texttt{║◌◌C◌◌┌─╢ 3} \\
\texttt{║◌◌◌┌─┤◌║ 4} \\
\texttt{║◌◌I│◌├─╢ 5} \\
\texttt{╟───┼─┘◌║ 6} \\
\texttt{║◌◌T│◌◌◌║ 7} \\
\texttt{╚═══╧═══╝} \\
\texttt{```} \\
\texttt{Trenutno igrate krog 5 od največ 20.} \\
\texttt{Skupaj ste zbrali 2 od 12 kazni.} \\
\texttt{Drugi igralec je premaknil predmet na svoji mreži.} \\
\\ 
\\ 
\texttt{Kakšen je vaš naslednji ukaz?} \\
            }
        }
    }
    & \\ \\

    \theutterance \stepcounter{utterance}  
    & \multicolumn{4}{p{0.6\linewidth}}{
        \cellcolor[rgb]{0.9,0.9,0.9}{
            \makecell[{{p{\linewidth}}}]{
                \texttt{\tiny{[P1$\rangle$GM]}}
                \texttt{PREMAKNI SE: I, (4, 4)} \\
            }
        }
    }
    & & \\ \\

    \theutterance \stepcounter{utterance}  
    & & & \multicolumn{2}{p{0.3\linewidth}}{
        \cellcolor[rgb]{0.9,0.9,0.9}{
            \makecell[{{p{\linewidth}}}]{
                \texttt{\tiny{[GM$|$GM]}}
                \texttt{PREMAKNI SE: I, (4, 4)} \\
            }
        }
    }
    & & \\ \\

    \theutterance \stepcounter{utterance}  
    & & & \multicolumn{2}{p{0.3\linewidth}}{
        \cellcolor[rgb]{0.9,0.9,0.9}{
            \makecell[{{p{\linewidth}}}]{
                \texttt{\tiny{[GM$|$GM]}}
                \texttt{Kazen: (4,4) ni prazen, vendar vsebuje '┌'.} \\
\\ 
\texttt{Skupaj ste zbrali 3 od 12 kazni.} \\
\texttt{Prosim, poskusite znova!} \\
            }
        }
    }
    & & \\ \\

    \theutterance \stepcounter{utterance}  
    & & \multicolumn{4}{p{0.6\linewidth}}{
        \cellcolor[rgb]{0.9,0.9,0.9}{
            \makecell[{{p{\linewidth}}}]{
                \texttt{\tiny{[P1$\langle$GM]}}
                \texttt{Kazen: (4,4) ni prazen, vendar vsebuje '┌'.} \\
\\ 
\texttt{Skupaj ste zbrali 3 od 12 kazni.} \\
\texttt{Prosim, poskusite znova!} \\
            }
        }
    }
    & \\ \\

    \theutterance \stepcounter{utterance}  
    & \multicolumn{4}{p{0.6\linewidth}}{
        \cellcolor[rgb]{0.9,0.9,0.9}{
            \makecell[{{p{\linewidth}}}]{
                \texttt{\tiny{[P1$\rangle$GM]}}
                \texttt{PREMAKNI SE: I, (5, 4)} \\
            }
        }
    }
    & & \\ \\

    \theutterance \stepcounter{utterance}  
    & & & \multicolumn{2}{p{0.3\linewidth}}{
        \cellcolor[rgb]{0.9,0.9,0.9}{
            \makecell[{{p{\linewidth}}}]{
                \texttt{\tiny{[GM$|$GM]}}
                \texttt{PREMAKNI SE: I, (5, 4)} \\
            }
        }
    }
    & & \\ \\

    \theutterance \stepcounter{utterance}  
    & & & \multicolumn{2}{p{0.3\linewidth}}{
        \cellcolor[rgb]{0.9,0.9,0.9}{
            \makecell[{{p{\linewidth}}}]{
                \texttt{\tiny{[GM$|$GM]}}
                \texttt{Kazen: (5,4) ni prazen, vendar vsebuje '─'.} \\
\\ 
\texttt{Skupaj ste zbrali 4 od 12 kazni.} \\
\texttt{Prosim, poskusite znova!} \\
            }
        }
    }
    & & \\ \\

    \theutterance \stepcounter{utterance}  
    & & \multicolumn{4}{p{0.6\linewidth}}{
        \cellcolor[rgb]{0.9,0.9,0.9}{
            \makecell[{{p{\linewidth}}}]{
                \texttt{\tiny{[P1$\langle$GM]}}
                \texttt{Kazen: (5,4) ni prazen, vendar vsebuje '─'.} \\
\\ 
\texttt{Skupaj ste zbrali 4 od 12 kazni.} \\
\texttt{Prosim, poskusite znova!} \\
            }
        }
    }
    & \\ \\

    \theutterance \stepcounter{utterance}  
    & \multicolumn{4}{p{0.6\linewidth}}{
        \cellcolor[rgb]{0.9,0.9,0.9}{
            \makecell[{{p{\linewidth}}}]{
                \texttt{\tiny{[P1$\rangle$GM]}}
                \texttt{PREMAKNI SE: I, (6, 4)} \\
            }
        }
    }
    & & \\ \\

    \theutterance \stepcounter{utterance}  
    & & & \multicolumn{2}{p{0.3\linewidth}}{
        \cellcolor[rgb]{0.9,0.9,0.9}{
            \makecell[{{p{\linewidth}}}]{
                \texttt{\tiny{[GM$|$GM]}}
                \texttt{PREMAKNI SE: I, (6, 4)} \\
            }
        }
    }
    & & \\ \\

    \theutterance \stepcounter{utterance}  
    & & & \multicolumn{2}{p{0.3\linewidth}}{
        \cellcolor[rgb]{0.9,0.9,0.9}{
            \makecell[{{p{\linewidth}}}]{
                \texttt{\tiny{[GM$|$GM]}}
                \texttt{Kazen: (6,4) ni prazen, vendar vsebuje '┤'.} \\
\\ 
\texttt{Skupaj ste zbrali 5 od 12 kazni.} \\
\texttt{Prosim, poskusite znova!} \\
            }
        }
    }
    & & \\ \\

    \theutterance \stepcounter{utterance}  
    & & \multicolumn{4}{p{0.6\linewidth}}{
        \cellcolor[rgb]{0.9,0.9,0.9}{
            \makecell[{{p{\linewidth}}}]{
                \texttt{\tiny{[P1$\langle$GM]}}
                \texttt{Kazen: (6,4) ni prazen, vendar vsebuje '┤'.} \\
\\ 
\texttt{Skupaj ste zbrali 5 od 12 kazni.} \\
\texttt{Prosim, poskusite znova!} \\
            }
        }
    }
    & \\ \\

    \theutterance \stepcounter{utterance}  
    & \multicolumn{4}{p{0.6\linewidth}}{
        \cellcolor[rgb]{0.9,0.9,0.9}{
            \makecell[{{p{\linewidth}}}]{
                \texttt{\tiny{[P1$\rangle$GM]}}
                \texttt{PREMAKNI SE: I, (7, 4)} \\
            }
        }
    }
    & & \\ \\

    \theutterance \stepcounter{utterance}  
    & & & \multicolumn{2}{p{0.3\linewidth}}{
        \cellcolor[rgb]{0.9,0.9,0.9}{
            \makecell[{{p{\linewidth}}}]{
                \texttt{\tiny{[GM$|$GM]}}
                \texttt{PREMAKNI SE: I, (7, 4)} \\
            }
        }
    }
    & & \\ \\

    \theutterance \stepcounter{utterance}  
    & & & \multicolumn{2}{p{0.3\linewidth}}{
        \cellcolor[rgb]{0.9,0.9,0.9}{
            \makecell[{{p{\linewidth}}}]{
                \texttt{\tiny{[GM$|$GM]}}
                \texttt{I moved to (7, 4)} \\
            }
        }
    }
    & & \\ \\

    \theutterance \stepcounter{utterance}  
    & & \multicolumn{4}{p{0.6\linewidth}}{
        \cellcolor[rgb]{0.9,0.9,0.9}{
            \makecell[{{p{\linewidth}}}]{
                \texttt{\tiny{[GM$\rangle$P2]}}
                \texttt{Uspešno premaknjen 'T' v (5,3). Vaša posodobljena mreža je videti takole:} \\
\texttt{```} \\
\texttt{ 1234567} \\
\texttt{╔═══════╗} \\
\texttt{║W◌◌◌◌◌◌║ 1} \\
\texttt{║◌H◌◌◌◌◌║ 2} \\
\texttt{║◌◌C◌T┌─╢ 3} \\
\texttt{║◌◌◌┌─┤◌║ 4} \\
\texttt{║◌◌◌│I├─╢ 5} \\
\texttt{╟───┼─┘◌║ 6} \\
\texttt{║◌◌◌│◌◌◌║ 7} \\
\texttt{╚═══╧═══╝} \\
\texttt{```} \\
\texttt{Trenutno igrate krog 6 od največ 20.} \\
\texttt{Skupaj ste zbrali 5 od 12 kazni.} \\
\texttt{Drugi igralec je premaknil predmet na svoji mreži.} \\
\\ 
\\ 
\texttt{Kakšen je vaš naslednji ukaz?} \\
            }
        }
    }
    & \\ \\

    \theutterance \stepcounter{utterance}  
    & & & \multicolumn{4}{p{0.6\linewidth}}{
        \cellcolor[rgb]{0.9,0.9,0.9}{
            \makecell[{{p{\linewidth}}}]{
                \texttt{\tiny{[GM$\langle$P2]}}
                \texttt{PREMAKNI SE: T, (5,5)} \\
            }
        }
    }
     \\ \\

    \theutterance \stepcounter{utterance}  
    & & & \multicolumn{2}{p{0.3\linewidth}}{
        \cellcolor[rgb]{0.9,0.9,0.9}{
            \makecell[{{p{\linewidth}}}]{
                \texttt{\tiny{[GM$|$GM]}}
                \texttt{PREMAKNI SE: T, (5,5)} \\
            }
        }
    }
    & & \\ \\

    \theutterance \stepcounter{utterance}  
    & & & \multicolumn{2}{p{0.3\linewidth}}{
        \cellcolor[rgb]{0.9,0.9,0.9}{
            \makecell[{{p{\linewidth}}}]{
                \texttt{\tiny{[GM$|$GM]}}
                \texttt{Kazen: (5,5) ni prazen, vendar vsebuje 'I'.} \\
\\ 
\texttt{Skupaj ste zbrali 6 od 12 kazni.} \\
\texttt{Prosim, poskusite znova!} \\
            }
        }
    }
    & & \\ \\

    \theutterance \stepcounter{utterance}  
    & & \multicolumn{4}{p{0.6\linewidth}}{
        \cellcolor[rgb]{0.9,0.9,0.9}{
            \makecell[{{p{\linewidth}}}]{
                \texttt{\tiny{[GM$\rangle$P2]}}
                \texttt{Kazen: (5,5) ni prazen, vendar vsebuje 'I'.} \\
\\ 
\texttt{Skupaj ste zbrali 6 od 12 kazni.} \\
\texttt{Prosim, poskusite znova!} \\
            }
        }
    }
    & \\ \\

    \theutterance \stepcounter{utterance}  
    & & & \multicolumn{4}{p{0.6\linewidth}}{
        \cellcolor[rgb]{0.9,0.9,0.9}{
            \makecell[{{p{\linewidth}}}]{
                \texttt{\tiny{[GM$\langle$P2]}}
                \texttt{PREMAKNI SE: I, (4,4)} \\
            }
        }
    }
     \\ \\

    \theutterance \stepcounter{utterance}  
    & & & \multicolumn{2}{p{0.3\linewidth}}{
        \cellcolor[rgb]{0.9,0.9,0.9}{
            \makecell[{{p{\linewidth}}}]{
                \texttt{\tiny{[GM$|$GM]}}
                \texttt{PREMAKNI SE: I, (4,4)} \\
            }
        }
    }
    & & \\ \\

    \theutterance \stepcounter{utterance}  
    & & & \multicolumn{2}{p{0.3\linewidth}}{
        \cellcolor[rgb]{0.9,0.9,0.9}{
            \makecell[{{p{\linewidth}}}]{
                \texttt{\tiny{[GM$|$GM]}}
                \texttt{Kazen: (4,4) ni prazen, vendar vsebuje '┌'.} \\
\\ 
\texttt{Skupaj ste zbrali 7 od 12 kazni.} \\
\texttt{Prosim, poskusite znova!} \\
            }
        }
    }
    & & \\ \\

    \theutterance \stepcounter{utterance}  
    & & \multicolumn{4}{p{0.6\linewidth}}{
        \cellcolor[rgb]{0.9,0.9,0.9}{
            \makecell[{{p{\linewidth}}}]{
                \texttt{\tiny{[GM$\rangle$P2]}}
                \texttt{Kazen: (4,4) ni prazen, vendar vsebuje '┌'.} \\
\\ 
\texttt{Skupaj ste zbrali 7 od 12 kazni.} \\
\texttt{Prosim, poskusite znova!} \\
            }
        }
    }
    & \\ \\

    \theutterance \stepcounter{utterance}  
    & & & \multicolumn{4}{p{0.6\linewidth}}{
        \cellcolor[rgb]{0.9,0.9,0.9}{
            \makecell[{{p{\linewidth}}}]{
                \texttt{\tiny{[GM$\langle$P2]}}
                \texttt{PREMAKNI SE: I, (5,5)} \\
            }
        }
    }
     \\ \\

    \theutterance \stepcounter{utterance}  
    & & & \multicolumn{2}{p{0.3\linewidth}}{
        \cellcolor[rgb]{0.9,0.9,0.9}{
            \makecell[{{p{\linewidth}}}]{
                \texttt{\tiny{[GM$|$GM]}}
                \texttt{PREMAKNI SE: I, (5,5)} \\
            }
        }
    }
    & & \\ \\

    \theutterance \stepcounter{utterance}  
    & & & \multicolumn{2}{p{0.3\linewidth}}{
        \cellcolor[rgb]{0.9,0.9,0.9}{
            \makecell[{{p{\linewidth}}}]{
                \texttt{\tiny{[GM$|$GM]}}
                \texttt{Kazen: (5,5) ni prazen, vendar vsebuje 'I'.} \\
\\ 
\texttt{Skupaj ste zbrali 8 od 12 kazni.} \\
\texttt{Prosim, poskusite znova!} \\
            }
        }
    }
    & & \\ \\

    \theutterance \stepcounter{utterance}  
    & & \multicolumn{4}{p{0.6\linewidth}}{
        \cellcolor[rgb]{0.9,0.9,0.9}{
            \makecell[{{p{\linewidth}}}]{
                \texttt{\tiny{[GM$\rangle$P2]}}
                \texttt{Kazen: (5,5) ni prazen, vendar vsebuje 'I'.} \\
\\ 
\texttt{Skupaj ste zbrali 8 od 12 kazni.} \\
\texttt{Prosim, poskusite znova!} \\
            }
        }
    }
    & \\ \\

    \theutterance \stepcounter{utterance}  
    & & & \multicolumn{4}{p{0.6\linewidth}}{
        \cellcolor[rgb]{0.9,0.9,0.9}{
            \makecell[{{p{\linewidth}}}]{
                \texttt{\tiny{[GM$\langle$P2]}}
                \texttt{RECI: Opazujem, da je I že na (5,5), kar je ciljna koordinata. Ne bom ga premikal. Pripravljam se za premik T na (5,5), a to zahteva, da (5,5) postane prosto. Ali lahko potrdite, ali je vaš I že na (4,4) ali drugje, da ne bi dvakrat zasedli istega cilja?} \\
            }
        }
    }
     \\ \\

    \theutterance \stepcounter{utterance}  
    & & & \multicolumn{2}{p{0.3\linewidth}}{
        \cellcolor[rgb]{0.9,0.9,0.9}{
            \makecell[{{p{\linewidth}}}]{
                \texttt{\tiny{[GM$|$GM]}}
                \texttt{RECI: Opazujem, da je I že na (5,5), kar je ciljna koordinata. Ne bom ga premikal. Pripravljam se za premik T na (5,5), a to zahteva, da (5,5) postane prosto. Ali lahko potrdite, ali je vaš I že na (4,4) ali drugje, da ne bi dvakrat zasedli istega cilja?} \\
            }
        }
    }
    & & \\ \\

    \theutterance \stepcounter{utterance}  
    & & \multicolumn{4}{p{0.6\linewidth}}{
        \cellcolor[rgb]{0.9,0.9,0.9}{
            \makecell[{{p{\linewidth}}}]{
                \texttt{\tiny{[P1$\langle$GM]}}
                \texttt{Uspešno premaknjen 'I' v (7,4). Vaša posodobljena mreža je videti takole:} \\
\texttt{```} \\
\texttt{ 1234567} \\
\texttt{╔═══════╗} \\
\texttt{║W◌◌◌◌◌◌║ 1} \\
\texttt{║◌H◌◌◌◌◌║ 2} \\
\texttt{║◌◌C◌◌┌─╢ 3} \\
\texttt{║◌◌◌┌─┤I║ 4} \\
\texttt{║◌◌◌│◌├─╢ 5} \\
\texttt{╟───┼─┘◌║ 6} \\
\texttt{║◌◌T│◌◌◌║ 7} \\
\texttt{╚═══╧═══╝} \\
\texttt{```} \\
\texttt{Trenutno igrate krog 6 od največ 20.} \\
\texttt{Skupaj ste zbrali 8 od 12 kazni.} \\
\texttt{Drugi igralec je poslal to sporočilo:} \\
\\ 
\\ 
\texttt{"Opazujem, da je I že na (5,5), kar je ciljna koordinata. Ne bom ga premikal. Pripravljam se za premik T na (5,5), a to zahteva, da (5,5) postane prosto. Ali lahko potrdite, ali je vaš I že na (4,4) ali drugje, da ne bi dvakrat zasedli istega cilja?"} \\
\\ 
\\ 
\texttt{Kakšen je vaš naslednji ukaz?} \\
            }
        }
    }
    & \\ \\

    \theutterance \stepcounter{utterance}  
    & \multicolumn{4}{p{0.6\linewidth}}{
        \cellcolor[rgb]{0.9,0.9,0.9}{
            \makecell[{{p{\linewidth}}}]{
                \texttt{\tiny{[P1$\rangle$GM]}}
                \texttt{RECI: Moj I je bil na (3,5) in ga sem premaknil na (7,4), ker je bila ciljna (4,4) zasedena s črto '┌'. Ciljna koordinata za I ostaja (4,4), a jo ne morem zasedati dokler je tam stena/premika. Kje je tvoj I? Če je na (5,5), je to drugačna koordinata kot naš dogovor (4,4). Ali naj I pustim na (7,4) ali poskušam najti drugo prosto celico? Tudi T je pri meni na (3,7).} \\
            }
        }
    }
    & & \\ \\

    \theutterance \stepcounter{utterance}  
    & & & \multicolumn{2}{p{0.3\linewidth}}{
        \cellcolor[rgb]{0.9,0.9,0.9}{
            \makecell[{{p{\linewidth}}}]{
                \texttt{\tiny{[GM$|$GM]}}
                \texttt{RECI: Moj I je bil na (3,5) in ga sem premaknil na (7,4), ker je bila ciljna (4,4) zasedena s črto '┌'. Ciljna koordinata za I ostaja (4,4), a jo ne morem zasedati dokler je tam stena/premika. Kje je tvoj I? Če je na (5,5), je to drugačna koordinata kot naš dogovor (4,4). Ali naj I pustim na (7,4) ali poskušam najti drugo prosto celico? Tudi T je pri meni na (3,7).} \\
            }
        }
    }
    & & \\ \\

    \theutterance \stepcounter{utterance}  
    & & \multicolumn{4}{p{0.6\linewidth}}{
        \cellcolor[rgb]{0.9,0.9,0.9}{
            \makecell[{{p{\linewidth}}}]{
                \texttt{\tiny{[GM$\rangle$P2]}}
                \texttt{Vaše sporočilo je bilo posredovano drugemu igralcu.} \\
\\ 
\\ 
\texttt{Trenutno igrate krog 7 od največ 20.} \\
\texttt{Skupaj ste zbrali 8 od 12 kazni.} \\
\texttt{Drugi igralec je poslal to sporočilo:} \\
\\ 
\\ 
\texttt{"Moj I je bil na (3,5) in ga sem premaknil na (7,4), ker je bila ciljna (4,4) zasedena s črto '┌'. Ciljna koordinata za I ostaja (4,4), a jo ne morem zasedati dokler je tam stena/premika. Kje je tvoj I? Če je na (5,5), je to drugačna koordinata kot naš dogovor (4,4). Ali naj I pustim na (7,4) ali poskušam najti drugo prosto celico? Tudi T je pri meni na (3,7)."} \\
\\ 
\\ 
\texttt{Kakšen je vaš naslednji ukaz?} \\
            }
        }
    }
    & \\ \\

    \theutterance \stepcounter{utterance}  
    & & & \multicolumn{4}{p{0.6\linewidth}}{
        \cellcolor[rgb]{0.9,0.9,0.9}{
            \makecell[{{p{\linewidth}}}]{
                \texttt{\tiny{[GM$\langle$P2]}}
                \texttt{RECI: Moj I je na (5,5). Ker je (4,4) zasedena s steno '┌', predlagam, da spremenimo cilj za I na (5,5) (tam je sedaj) in za T na (4,5) (tam je sedaj prazno, ko gre I stran). To preprečuje več kazni. Strinjate se?} \\
            }
        }
    }
     \\ \\

    \theutterance \stepcounter{utterance}  
    & & & \multicolumn{2}{p{0.3\linewidth}}{
        \cellcolor[rgb]{0.9,0.9,0.9}{
            \makecell[{{p{\linewidth}}}]{
                \texttt{\tiny{[GM$|$GM]}}
                \texttt{RECI: Moj I je na (5,5). Ker je (4,4) zasedena s steno '┌', predlagam, da spremenimo cilj za I na (5,5) (tam je sedaj) in za T na (4,5) (tam je sedaj prazno, ko gre I stran). To preprečuje več kazni. Strinjate se?} \\
            }
        }
    }
    & & \\ \\

    \theutterance \stepcounter{utterance}  
    & & \multicolumn{4}{p{0.6\linewidth}}{
        \cellcolor[rgb]{0.9,0.9,0.9}{
            \makecell[{{p{\linewidth}}}]{
                \texttt{\tiny{[P1$\langle$GM]}}
                \texttt{Vaše sporočilo je bilo posredovano drugemu igralcu.} \\
\\ 
\\ 
\texttt{Trenutno igrate krog 7 od največ 20.} \\
\texttt{Skupaj ste zbrali 8 od 12 kazni.} \\
\texttt{Drugi igralec je poslal to sporočilo:} \\
\\ 
\\ 
\texttt{"Moj I je na (5,5). Ker je (4,4) zasedena s steno '┌', predlagam, da spremenimo cilj za I na (5,5) (tam je sedaj) in za T na (4,5) (tam je sedaj prazno, ko gre I stran). To preprečuje več kazni. Strinjate se?"} \\
\\ 
\\ 
\texttt{Kakšen je vaš naslednji ukaz?} \\
            }
        }
    }
    & \\ \\

    \theutterance \stepcounter{utterance}  
    & \multicolumn{4}{p{0.6\linewidth}}{
        \cellcolor[rgb]{0.9,0.9,0.9}{
            \makecell[{{p{\linewidth}}}]{
                \texttt{\tiny{[P1$\rangle$GM]}}
                \texttt{RECI: Ne strinjam se s spremembo cilja za I na (5,5), ker je ta celica zasedena s tvojim I (ne more biti hkrati na obeh ploščah na isti koordinati, če je že zasedena). Predlagam, da I pustimo na (7,4) ali najdemo drugo prosto ciljno koordinato, npr. (1,3) ali (1,4), ki sta prosta pri meni. Za T pa predlagam cilj (4,5), ki je pri meni prosto. Ali strinjam, da T gre na (4,5)? Kje je tvoj T sedaj?} \\
            }
        }
    }
    & & \\ \\

    \theutterance \stepcounter{utterance}  
    & & & \multicolumn{2}{p{0.3\linewidth}}{
        \cellcolor[rgb]{0.9,0.9,0.9}{
            \makecell[{{p{\linewidth}}}]{
                \texttt{\tiny{[GM$|$GM]}}
                \texttt{RECI: Ne strinjam se s spremembo cilja za I na (5,5), ker je ta celica zasedena s tvojim I (ne more biti hkrati na obeh ploščah na isti koordinati, če je že zasedena). Predlagam, da I pustimo na (7,4) ali najdemo drugo prosto ciljno koordinato, npr. (1,3) ali (1,4), ki sta prosta pri meni. Za T pa predlagam cilj (4,5), ki je pri meni prosto. Ali strinjam, da T gre na (4,5)? Kje je tvoj T sedaj?} \\
            }
        }
    }
    & & \\ \\

    \theutterance \stepcounter{utterance}  
    & & \multicolumn{4}{p{0.6\linewidth}}{
        \cellcolor[rgb]{0.9,0.9,0.9}{
            \makecell[{{p{\linewidth}}}]{
                \texttt{\tiny{[GM$\rangle$P2]}}
                \texttt{Vaše sporočilo je bilo posredovano drugemu igralcu.} \\
\\ 
\\ 
\texttt{Trenutno igrate krog 8 od največ 20.} \\
\texttt{Skupaj ste zbrali 8 od 12 kazni.} \\
\texttt{Drugi igralec je poslal to sporočilo:} \\
\\ 
\\ 
\texttt{"Ne strinjam se s spremembo cilja za I na (5,5), ker je ta celica zasedena s tvojim I (ne more biti hkrati na obeh ploščah na isti koordinati, če je že zasedena). Predlagam, da I pustimo na (7,4) ali najdemo drugo prosto ciljno koordinato, npr. (1,3) ali (1,4), ki sta prosta pri meni. Za T pa predlagam cilj (4,5), ki je pri meni prosto. Ali strinjam, da T gre na (4,5)? Kje je tvoj T sedaj?"} \\
\\ 
\\ 
\texttt{Kakšen je vaš naslednji ukaz?} \\
            }
        }
    }
    & \\ \\

    \theutterance \stepcounter{utterance}  
    & & & \multicolumn{4}{p{0.6\linewidth}}{
        \cellcolor[rgb]{0.9,0.9,0.9}{
            \makecell[{{p{\linewidth}}}]{
                \texttt{\tiny{[GM$\langle$P2]}}
                \texttt{RECI: Moj T je na (5,3). Strinjam se, da T gre na (4,5). Za I predlagam cilj (6,5), ki je pri meni prosto. Ali je (6,5) prosto tudi pri tebi?} \\
            }
        }
    }
     \\ \\

    \theutterance \stepcounter{utterance}  
    & & & \multicolumn{2}{p{0.3\linewidth}}{
        \cellcolor[rgb]{0.9,0.9,0.9}{
            \makecell[{{p{\linewidth}}}]{
                \texttt{\tiny{[GM$|$GM]}}
                \texttt{RECI: Moj T je na (5,3). Strinjam se, da T gre na (4,5). Za I predlagam cilj (6,5), ki je pri meni prosto. Ali je (6,5) prosto tudi pri tebi?} \\
            }
        }
    }
    & & \\ \\

    \theutterance \stepcounter{utterance}  
    & & \multicolumn{4}{p{0.6\linewidth}}{
        \cellcolor[rgb]{0.9,0.9,0.9}{
            \makecell[{{p{\linewidth}}}]{
                \texttt{\tiny{[P1$\langle$GM]}}
                \texttt{Vaše sporočilo je bilo posredovano drugemu igralcu.} \\
\\ 
\\ 
\texttt{Trenutno igrate krog 8 od največ 20.} \\
\texttt{Skupaj ste zbrali 8 od 12 kazni.} \\
\texttt{Drugi igralec je poslal to sporočilo:} \\
\\ 
\\ 
\texttt{"Moj T je na (5,3). Strinjam se, da T gre na (4,5). Za I predlagam cilj (6,5), ki je pri meni prosto. Ali je (6,5) prosto tudi pri tebi?"} \\
\\ 
\\ 
\texttt{Kakšen je vaš naslednji ukaz?} \\
            }
        }
    }
    & \\ \\

    \theutterance \stepcounter{utterance}  
    & \multicolumn{4}{p{0.6\linewidth}}{
        \cellcolor[rgb]{0.9,0.9,0.9}{
            \makecell[{{p{\linewidth}}}]{
                \texttt{\tiny{[P1$\rangle$GM]}}
                \texttt{PREMAKNI SE: T, (4, 5)} \\
            }
        }
    }
    & & \\ \\

    \theutterance \stepcounter{utterance}  
    & & & \multicolumn{2}{p{0.3\linewidth}}{
        \cellcolor[rgb]{0.9,0.9,0.9}{
            \makecell[{{p{\linewidth}}}]{
                \texttt{\tiny{[GM$|$GM]}}
                \texttt{PREMAKNI SE: T, (4, 5)} \\
            }
        }
    }
    & & \\ \\

    \theutterance \stepcounter{utterance}  
    & & & \multicolumn{2}{p{0.3\linewidth}}{
        \cellcolor[rgb]{0.9,0.9,0.9}{
            \makecell[{{p{\linewidth}}}]{
                \texttt{\tiny{[GM$|$GM]}}
                \texttt{Kazen: (4,5) ni prazen, vendar vsebuje '│'.} \\
\\ 
\texttt{Skupaj ste zbrali 9 od 12 kazni.} \\
\texttt{Prosim, poskusite znova!} \\
            }
        }
    }
    & & \\ \\

    \theutterance \stepcounter{utterance}  
    & & \multicolumn{4}{p{0.6\linewidth}}{
        \cellcolor[rgb]{0.9,0.9,0.9}{
            \makecell[{{p{\linewidth}}}]{
                \texttt{\tiny{[P1$\langle$GM]}}
                \texttt{Kazen: (4,5) ni prazen, vendar vsebuje '│'.} \\
\\ 
\texttt{Skupaj ste zbrali 9 od 12 kazni.} \\
\texttt{Prosim, poskusite znova!} \\
            }
        }
    }
    & \\ \\

    \theutterance \stepcounter{utterance}  
    & \multicolumn{4}{p{0.6\linewidth}}{
        \cellcolor[rgb]{0.9,0.9,0.9}{
            \makecell[{{p{\linewidth}}}]{
                \texttt{\tiny{[P1$\rangle$GM]}}
                \texttt{RECI: (4,5) je pri meni zasedeno s '│'. Poskušam drug cilj za T. Ali je (5,4) prosto pri tebi? Pri meni je (5,4) prosto. Ali (6,5) za I? Pri meni je (6,5) prosto. Predlagam T {-}> (5,4), I {-}> (6,5). Strinjam se?} \\
            }
        }
    }
    & & \\ \\

    \theutterance \stepcounter{utterance}  
    & & & \multicolumn{2}{p{0.3\linewidth}}{
        \cellcolor[rgb]{0.9,0.9,0.9}{
            \makecell[{{p{\linewidth}}}]{
                \texttt{\tiny{[GM$|$GM]}}
                \texttt{RECI: (4,5) je pri meni zasedeno s '│'. Poskušam drug cilj za T. Ali je (5,4) prosto pri tebi? Pri meni je (5,4) prosto. Ali (6,5) za I? Pri meni je (6,5) prosto. Predlagam T {-}> (5,4), I {-}> (6,5). Strinjam se?} \\
            }
        }
    }
    & & \\ \\

    \theutterance \stepcounter{utterance}  
    & & \multicolumn{4}{p{0.6\linewidth}}{
        \cellcolor[rgb]{0.9,0.9,0.9}{
            \makecell[{{p{\linewidth}}}]{
                \texttt{\tiny{[GM$\rangle$P2]}}
                \texttt{Vaše sporočilo je bilo posredovano drugemu igralcu.} \\
\\ 
\\ 
\texttt{Trenutno igrate krog 9 od največ 20.} \\
\texttt{Skupaj ste zbrali 9 od 12 kazni.} \\
\texttt{Drugi igralec je poslal to sporočilo:} \\
\\ 
\\ 
\texttt{"(4,5) je pri meni zasedeno s '│'. Poskušam drug cilj za T. Ali je (5,4) prosto pri tebi? Pri meni je (5,4) prosto. Ali (6,5) za I? Pri meni je (6,5) prosto. Predlagam T {-}> (5,4), I {-}> (6,5). Strinjam se?"} \\
\\ 
\\ 
\texttt{Kakšen je vaš naslednji ukaz?} \\
            }
        }
    }
    & \\ \\

    \theutterance \stepcounter{utterance}  
    & & & \multicolumn{4}{p{0.6\linewidth}}{
        \cellcolor[rgb]{0.9,0.9,0.9}{
            \makecell[{{p{\linewidth}}}]{
                \texttt{\tiny{[GM$\langle$P2]}}
                \texttt{RECI: Pri meni je (5,4) zasedena s '┌', (6,5) pa je prosto. Predlagam T {-}> (6,4) (prosto pri meni), I {-}> (6,5) (prosto pri meni). Preveri, ali sta oba prosta tudi pri tebi.} \\
            }
        }
    }
     \\ \\

    \theutterance \stepcounter{utterance}  
    & & & \multicolumn{2}{p{0.3\linewidth}}{
        \cellcolor[rgb]{0.9,0.9,0.9}{
            \makecell[{{p{\linewidth}}}]{
                \texttt{\tiny{[GM$|$GM]}}
                \texttt{RECI: Pri meni je (5,4) zasedena s '┌', (6,5) pa je prosto. Predlagam T {-}> (6,4) (prosto pri meni), I {-}> (6,5) (prosto pri meni). Preveri, ali sta oba prosta tudi pri tebi.} \\
            }
        }
    }
    & & \\ \\

    \theutterance \stepcounter{utterance}  
    & & \multicolumn{4}{p{0.6\linewidth}}{
        \cellcolor[rgb]{0.9,0.9,0.9}{
            \makecell[{{p{\linewidth}}}]{
                \texttt{\tiny{[P1$\langle$GM]}}
                \texttt{Vaše sporočilo je bilo posredovano drugemu igralcu.} \\
\\ 
\\ 
\texttt{Trenutno igrate krog 9 od največ 20.} \\
\texttt{Skupaj ste zbrali 9 od 12 kazni.} \\
\texttt{Drugi igralec je poslal to sporočilo:} \\
\\ 
\\ 
\texttt{"Pri meni je (5,4) zasedena s '┌', (6,5) pa je prosto. Predlagam T {-}> (6,4) (prosto pri meni), I {-}> (6,5) (prosto pri meni). Preveri, ali sta oba prosta tudi pri tebi."} \\
\\ 
\\ 
\texttt{Kakšen je vaš naslednji ukaz?} \\
            }
        }
    }
    & \\ \\

    \theutterance \stepcounter{utterance}  
    & \multicolumn{4}{p{0.6\linewidth}}{
        \cellcolor[rgb]{0.9,0.9,0.9}{
            \makecell[{{p{\linewidth}}}]{
                \texttt{\tiny{[P1$\rangle$GM]}}
                \texttt{PREMAKNI SE: I, (6, 5)} \\
            }
        }
    }
    & & \\ \\

    \theutterance \stepcounter{utterance}  
    & & & \multicolumn{2}{p{0.3\linewidth}}{
        \cellcolor[rgb]{0.9,0.9,0.9}{
            \makecell[{{p{\linewidth}}}]{
                \texttt{\tiny{[GM$|$GM]}}
                \texttt{PREMAKNI SE: I, (6, 5)} \\
            }
        }
    }
    & & \\ \\

    \theutterance \stepcounter{utterance}  
    & & & \multicolumn{2}{p{0.3\linewidth}}{
        \cellcolor[rgb]{0.9,0.9,0.9}{
            \makecell[{{p{\linewidth}}}]{
                \texttt{\tiny{[GM$|$GM]}}
                \texttt{Kazen: (6,5) ni prazen, vendar vsebuje '├'.} \\
\\ 
\texttt{Skupaj ste zbrali 10 od 12 kazni.} \\
\texttt{Prosim, poskusite znova!} \\
            }
        }
    }
    & & \\ \\

    \theutterance \stepcounter{utterance}  
    & & \multicolumn{4}{p{0.6\linewidth}}{
        \cellcolor[rgb]{0.9,0.9,0.9}{
            \makecell[{{p{\linewidth}}}]{
                \texttt{\tiny{[P1$\langle$GM]}}
                \texttt{Kazen: (6,5) ni prazen, vendar vsebuje '├'.} \\
\\ 
\texttt{Skupaj ste zbrali 10 od 12 kazni.} \\
\texttt{Prosim, poskusite znova!} \\
            }
        }
    }
    & \\ \\

    \theutterance \stepcounter{utterance}  
    & \multicolumn{4}{p{0.6\linewidth}}{
        \cellcolor[rgb]{0.9,0.9,0.9}{
            \makecell[{{p{\linewidth}}}]{
                \texttt{\tiny{[P1$\rangle$GM]}}
                \texttt{PREMAKNI SE: T, (6, 4)} \\
            }
        }
    }
    & & \\ \\

    \theutterance \stepcounter{utterance}  
    & & & \multicolumn{2}{p{0.3\linewidth}}{
        \cellcolor[rgb]{0.9,0.9,0.9}{
            \makecell[{{p{\linewidth}}}]{
                \texttt{\tiny{[GM$|$GM]}}
                \texttt{PREMAKNI SE: T, (6, 4)} \\
            }
        }
    }
    & & \\ \\

    \theutterance \stepcounter{utterance}  
    & & & \multicolumn{2}{p{0.3\linewidth}}{
        \cellcolor[rgb]{0.9,0.9,0.9}{
            \makecell[{{p{\linewidth}}}]{
                \texttt{\tiny{[GM$|$GM]}}
                \texttt{Kazen: (6,4) ni prazen, vendar vsebuje '┤'.} \\
\\ 
\texttt{Skupaj ste zbrali 11 od 12 kazni.} \\
\texttt{Prosim, poskusite znova!} \\
            }
        }
    }
    & & \\ \\

    \theutterance \stepcounter{utterance}  
    & & \multicolumn{4}{p{0.6\linewidth}}{
        \cellcolor[rgb]{0.9,0.9,0.9}{
            \makecell[{{p{\linewidth}}}]{
                \texttt{\tiny{[P1$\langle$GM]}}
                \texttt{Kazen: (6,4) ni prazen, vendar vsebuje '┤'.} \\
\\ 
\texttt{Skupaj ste zbrali 11 od 12 kazni.} \\
\texttt{Prosim, poskusite znova!} \\
            }
        }
    }
    & \\ \\

    \theutterance \stepcounter{utterance}  
    & \multicolumn{4}{p{0.6\linewidth}}{
        \cellcolor[rgb]{0.9,0.9,0.9}{
            \makecell[{{p{\linewidth}}}]{
                \texttt{\tiny{[P1$\rangle$GM]}}
                \texttt{RECI: Pri meni so (6,4) in (6,5) zasedeni s stena ('┤' in '├'). Prosti so (1,3), (1,4), (7,3), (7,5), (7,6), (7,7), (5,4). Ali kje imaš prosto za T in I? Pozdravljen!} \\
            }
        }
    }
    & & \\ \\

    \theutterance \stepcounter{utterance}  
    & & & \multicolumn{2}{p{0.3\linewidth}}{
        \cellcolor[rgb]{0.9,0.9,0.9}{
            \makecell[{{p{\linewidth}}}]{
                \texttt{\tiny{[GM$|$GM]}}
                \texttt{RECI: Pri meni so (6,4) in (6,5) zasedeni s stena ('┤' in '├'). Prosti so (1,3), (1,4), (7,3), (7,5), (7,6), (7,7), (5,4). Ali kje imaš prosto za T in I? Pozdravljen!} \\
            }
        }
    }
    & & \\ \\

    \theutterance \stepcounter{utterance}  
    & & \multicolumn{4}{p{0.6\linewidth}}{
        \cellcolor[rgb]{0.9,0.9,0.9}{
            \makecell[{{p{\linewidth}}}]{
                \texttt{\tiny{[GM$\rangle$P2]}}
                \texttt{Vaše sporočilo je bilo posredovano drugemu igralcu.} \\
\\ 
\\ 
\texttt{Trenutno igrate krog 10 od največ 20.} \\
\texttt{Skupaj ste zbrali 11 od 12 kazni.} \\
\texttt{Drugi igralec je poslal to sporočilo:} \\
\\ 
\\ 
\texttt{"Pri meni so (6,4) in (6,5) zasedeni s stena ('┤' in '├'). Prosti so (1,3), (1,4), (7,3), (7,5), (7,6), (7,7), (5,4). Ali kje imaš prosto za T in I? Pozdravljen!"} \\
\\ 
\\ 
\texttt{Kakšen je vaš naslednji ukaz?} \\
            }
        }
    }
    & \\ \\

    \theutterance \stepcounter{utterance}  
    & & & \multicolumn{4}{p{0.6\linewidth}}{
        \cellcolor[rgb]{0.9,0.9,0.9}{
            \makecell[{{p{\linewidth}}}]{
                \texttt{\tiny{[GM$\langle$P2]}}
                \texttt{RECI: Pri meni so proste: (1,3), (1,4), (7,3), (7,4), (7,5), (7,6), (7,7). Skupne proste celice: (1,3), (1,4), (7,3), (7,4), (7,5), (7,6), (7,7). Predlagam T {-}> (1,3), I {-}> (1,4). Strinjam se?} \\
            }
        }
    }
     \\ \\

    \theutterance \stepcounter{utterance}  
    & & & \multicolumn{2}{p{0.3\linewidth}}{
        \cellcolor[rgb]{0.9,0.9,0.9}{
            \makecell[{{p{\linewidth}}}]{
                \texttt{\tiny{[GM$|$GM]}}
                \texttt{RECI: Pri meni so proste: (1,3), (1,4), (7,3), (7,4), (7,5), (7,6), (7,7). Skupne proste celice: (1,3), (1,4), (7,3), (7,4), (7,5), (7,6), (7,7). Predlagam T {-}> (1,3), I {-}> (1,4). Strinjam se?} \\
            }
        }
    }
    & & \\ \\

    \theutterance \stepcounter{utterance}  
    & & \multicolumn{4}{p{0.6\linewidth}}{
        \cellcolor[rgb]{0.9,0.9,0.9}{
            \makecell[{{p{\linewidth}}}]{
                \texttt{\tiny{[P1$\langle$GM]}}
                \texttt{Vaše sporočilo je bilo posredovano drugemu igralcu.} \\
\\ 
\\ 
\texttt{Trenutno igrate krog 10 od največ 20.} \\
\texttt{Skupaj ste zbrali 11 od 12 kazni.} \\
\texttt{Drugi igralec je poslal to sporočilo:} \\
\\ 
\\ 
\texttt{"Pri meni so proste: (1,3), (1,4), (7,3), (7,4), (7,5), (7,6), (7,7). Skupne proste celice: (1,3), (1,4), (7,3), (7,4), (7,5), (7,6), (7,7). Predlagam T {-}> (1,3), I {-}> (1,4). Strinjam se?"} \\
\\ 
\\ 
\texttt{Kakšen je vaš naslednji ukaz?} \\
            }
        }
    }
    & \\ \\

    \theutterance \stepcounter{utterance}  
    & \multicolumn{4}{p{0.6\linewidth}}{
        \cellcolor[rgb]{0.9,0.9,0.9}{
            \makecell[{{p{\linewidth}}}]{
                \texttt{\tiny{[P1$\rangle$GM]}}
                \texttt{RECI: Strinjam se: T {-}> (1,3), I {-}> (1,4). Izvedam premike.} \\
            }
        }
    }
    & & \\ \\

    \theutterance \stepcounter{utterance}  
    & & & \multicolumn{2}{p{0.3\linewidth}}{
        \cellcolor[rgb]{0.9,0.9,0.9}{
            \makecell[{{p{\linewidth}}}]{
                \texttt{\tiny{[GM$|$GM]}}
                \texttt{RECI: Strinjam se: T {-}> (1,3), I {-}> (1,4). Izvedam premike.} \\
            }
        }
    }
    & & \\ \\

    \theutterance \stepcounter{utterance}  
    & & \multicolumn{4}{p{0.6\linewidth}}{
        \cellcolor[rgb]{0.9,0.9,0.9}{
            \makecell[{{p{\linewidth}}}]{
                \texttt{\tiny{[GM$\rangle$P2]}}
                \texttt{Vaše sporočilo je bilo posredovano drugemu igralcu.} \\
\\ 
\\ 
\texttt{Trenutno igrate krog 11 od največ 20.} \\
\texttt{Skupaj ste zbrali 11 od 12 kazni.} \\
\texttt{Drugi igralec je poslal to sporočilo:} \\
\\ 
\\ 
\texttt{"Strinjam se: T {-}> (1,3), I {-}> (1,4). Izvedam premike."} \\
\\ 
\\ 
\texttt{Kakšen je vaš naslednji ukaz?} \\
            }
        }
    }
    & \\ \\

    \theutterance \stepcounter{utterance}  
    & & & \multicolumn{4}{p{0.6\linewidth}}{
        \cellcolor[rgb]{0.9,0.9,0.9}{
            \makecell[{{p{\linewidth}}}]{
                \texttt{\tiny{[GM$\langle$P2]}}
                \texttt{PREMAKNI SE: T, (1,3)} \\
            }
        }
    }
     \\ \\

    \theutterance \stepcounter{utterance}  
    & & & \multicolumn{2}{p{0.3\linewidth}}{
        \cellcolor[rgb]{0.9,0.9,0.9}{
            \makecell[{{p{\linewidth}}}]{
                \texttt{\tiny{[GM$|$GM]}}
                \texttt{PREMAKNI SE: T, (1,3)} \\
            }
        }
    }
    & & \\ \\

    \theutterance \stepcounter{utterance}  
    & & & \multicolumn{2}{p{0.3\linewidth}}{
        \cellcolor[rgb]{0.9,0.9,0.9}{
            \makecell[{{p{\linewidth}}}]{
                \texttt{\tiny{[GM$|$GM]}}
                \texttt{T moved to (1, 3)} \\
            }
        }
    }
    & & \\ \\

    \theutterance \stepcounter{utterance}  
    & & \multicolumn{4}{p{0.6\linewidth}}{
        \cellcolor[rgb]{0.9,0.9,0.9}{
            \makecell[{{p{\linewidth}}}]{
                \texttt{\tiny{[P1$\langle$GM]}}
                \texttt{Vaše sporočilo je bilo posredovano drugemu igralcu.} \\
\\ 
\\ 
\texttt{Trenutno igrate krog 11 od največ 20.} \\
\texttt{Skupaj ste zbrali 11 od 12 kazni.} \\
\texttt{Drugi igralec je premaknil predmet na svoji mreži.} \\
\\ 
\\ 
\texttt{Kakšen je vaš naslednji ukaz?} \\
            }
        }
    }
    & \\ \\

    \theutterance \stepcounter{utterance}  
    & \multicolumn{4}{p{0.6\linewidth}}{
        \cellcolor[rgb]{0.9,0.9,0.9}{
            \makecell[{{p{\linewidth}}}]{
                \texttt{\tiny{[P1$\rangle$GM]}}
                \texttt{PREMAKNI SE: T, (1, 3)} \\
            }
        }
    }
    & & \\ \\

    \theutterance \stepcounter{utterance}  
    & & & \multicolumn{2}{p{0.3\linewidth}}{
        \cellcolor[rgb]{0.9,0.9,0.9}{
            \makecell[{{p{\linewidth}}}]{
                \texttt{\tiny{[GM$|$GM]}}
                \texttt{PREMAKNI SE: T, (1, 3)} \\
            }
        }
    }
    & & \\ \\

    \theutterance \stepcounter{utterance}  
    & & & \multicolumn{2}{p{0.3\linewidth}}{
        \cellcolor[rgb]{0.9,0.9,0.9}{
            \makecell[{{p{\linewidth}}}]{
                \texttt{\tiny{[GM$|$GM]}}
                \texttt{T moved to (1, 3)} \\
            }
        }
    }
    & & \\ \\

    \theutterance \stepcounter{utterance}  
    & & \multicolumn{4}{p{0.6\linewidth}}{
        \cellcolor[rgb]{0.9,0.9,0.9}{
            \makecell[{{p{\linewidth}}}]{
                \texttt{\tiny{[GM$\rangle$P2]}}
                \texttt{Uspešno premaknjen 'T' v (1,3). Vaša posodobljena mreža je videti takole:} \\
\texttt{```} \\
\texttt{ 1234567} \\
\texttt{╔═══════╗} \\
\texttt{║W◌◌◌◌◌◌║ 1} \\
\texttt{║◌H◌◌◌◌◌║ 2} \\
\texttt{║T◌C◌◌┌─╢ 3} \\
\texttt{║◌◌◌┌─┤◌║ 4} \\
\texttt{║◌◌◌│I├─╢ 5} \\
\texttt{╟───┼─┘◌║ 6} \\
\texttt{║◌◌◌│◌◌◌║ 7} \\
\texttt{╚═══╧═══╝} \\
\texttt{```} \\
\texttt{Trenutno igrate krog 12 od največ 20.} \\
\texttt{Skupaj ste zbrali 11 od 12 kazni.} \\
\texttt{Drugi igralec je premaknil predmet na svoji mreži.} \\
\\ 
\\ 
\texttt{Kakšen je vaš naslednji ukaz?} \\
            }
        }
    }
    & \\ \\

    \theutterance \stepcounter{utterance}  
    & & & \multicolumn{4}{p{0.6\linewidth}}{
        \cellcolor[rgb]{0.9,0.9,0.9}{
            \makecell[{{p{\linewidth}}}]{
                \texttt{\tiny{[GM$\langle$P2]}}
                \texttt{PREMAKNI SE: I, (1,4)} \\
            }
        }
    }
     \\ \\

    \theutterance \stepcounter{utterance}  
    & & & \multicolumn{2}{p{0.3\linewidth}}{
        \cellcolor[rgb]{0.9,0.9,0.9}{
            \makecell[{{p{\linewidth}}}]{
                \texttt{\tiny{[GM$|$GM]}}
                \texttt{PREMAKNI SE: I, (1,4)} \\
            }
        }
    }
    & & \\ \\

    \theutterance \stepcounter{utterance}  
    & & & \multicolumn{2}{p{0.3\linewidth}}{
        \cellcolor[rgb]{0.9,0.9,0.9}{
            \makecell[{{p{\linewidth}}}]{
                \texttt{\tiny{[GM$|$GM]}}
                \texttt{I moved to (1, 4)} \\
            }
        }
    }
    & & \\ \\

    \theutterance \stepcounter{utterance}  
    & & \multicolumn{4}{p{0.6\linewidth}}{
        \cellcolor[rgb]{0.9,0.9,0.9}{
            \makecell[{{p{\linewidth}}}]{
                \texttt{\tiny{[P1$\langle$GM]}}
                \texttt{Uspešno premaknjen 'T' v (1,3). Vaša posodobljena mreža je videti takole:} \\
\texttt{```} \\
\texttt{ 1234567} \\
\texttt{╔═══════╗} \\
\texttt{║W◌◌◌◌◌◌║ 1} \\
\texttt{║◌H◌◌◌◌◌║ 2} \\
\texttt{║T◌C◌◌┌─╢ 3} \\
\texttt{║◌◌◌┌─┤I║ 4} \\
\texttt{║◌◌◌│◌├─╢ 5} \\
\texttt{╟───┼─┘◌║ 6} \\
\texttt{║◌◌◌│◌◌◌║ 7} \\
\texttt{╚═══╧═══╝} \\
\texttt{```} \\
\texttt{Trenutno igrate krog 12 od največ 20.} \\
\texttt{Skupaj ste zbrali 11 od 12 kazni.} \\
\texttt{Drugi igralec je premaknil predmet na svoji mreži.} \\
\\ 
\\ 
\texttt{Kakšen je vaš naslednji ukaz?} \\
            }
        }
    }
    & \\ \\

    \theutterance \stepcounter{utterance}  
    & \multicolumn{4}{p{0.6\linewidth}}{
        \cellcolor[rgb]{0.9,0.9,0.9}{
            \makecell[{{p{\linewidth}}}]{
                \texttt{\tiny{[P1$\rangle$GM]}}
                \texttt{PREMAKNI SE: I, (1, 4)} \\
            }
        }
    }
    & & \\ \\

    \theutterance \stepcounter{utterance}  
    & & & \multicolumn{2}{p{0.3\linewidth}}{
        \cellcolor[rgb]{0.9,0.9,0.9}{
            \makecell[{{p{\linewidth}}}]{
                \texttt{\tiny{[GM$|$GM]}}
                \texttt{PREMAKNI SE: I, (1, 4)} \\
            }
        }
    }
    & & \\ \\

    \theutterance \stepcounter{utterance}  
    & & & \multicolumn{2}{p{0.3\linewidth}}{
        \cellcolor[rgb]{0.9,0.9,0.9}{
            \makecell[{{p{\linewidth}}}]{
                \texttt{\tiny{[GM$|$GM]}}
                \texttt{I moved to (1, 4)} \\
            }
        }
    }
    & & \\ \\

    \theutterance \stepcounter{utterance}  
    & & \multicolumn{4}{p{0.6\linewidth}}{
        \cellcolor[rgb]{0.9,0.9,0.9}{
            \makecell[{{p{\linewidth}}}]{
                \texttt{\tiny{[GM$\rangle$P2]}}
                \texttt{Uspešno premaknjen 'I' v (1,4). Vaša posodobljena mreža je videti takole:} \\
\texttt{```} \\
\texttt{ 1234567} \\
\texttt{╔═══════╗} \\
\texttt{║W◌◌◌◌◌◌║ 1} \\
\texttt{║◌H◌◌◌◌◌║ 2} \\
\texttt{║T◌C◌◌┌─╢ 3} \\
\texttt{║I◌◌┌─┤◌║ 4} \\
\texttt{║◌◌◌│◌├─╢ 5} \\
\texttt{╟───┼─┘◌║ 6} \\
\texttt{║◌◌◌│◌◌◌║ 7} \\
\texttt{╚═══╧═══╝} \\
\texttt{```} \\
\texttt{Trenutno igrate krog 13 od največ 20.} \\
\texttt{Skupaj ste zbrali 11 od 12 kazni.} \\
\texttt{Drugi igralec je premaknil predmet na svoji mreži.} \\
\\ 
\\ 
\texttt{Kakšen je vaš naslednji ukaz?} \\
            }
        }
    }
    & \\ \\

    \theutterance \stepcounter{utterance}  
    & & & \multicolumn{4}{p{0.6\linewidth}}{
        \cellcolor[rgb]{0.9,0.9,0.9}{
            \makecell[{{p{\linewidth}}}]{
                \texttt{\tiny{[GM$\langle$P2]}}
                \texttt{RECI: končano?} \\
            }
        }
    }
     \\ \\

    \theutterance \stepcounter{utterance}  
    & & & \multicolumn{2}{p{0.3\linewidth}}{
        \cellcolor[rgb]{0.9,0.9,0.9}{
            \makecell[{{p{\linewidth}}}]{
                \texttt{\tiny{[GM$|$GM]}}
                \texttt{RECI: končano?} \\
            }
        }
    }
    & & \\ \\

    \theutterance \stepcounter{utterance}  
    & & \multicolumn{4}{p{0.6\linewidth}}{
        \cellcolor[rgb]{0.9,0.9,0.9}{
            \makecell[{{p{\linewidth}}}]{
                \texttt{\tiny{[P1$\langle$GM]}}
                \texttt{Uspešno premaknjen 'I' v (1,4). Vaša posodobljena mreža je videti takole:} \\
\texttt{```} \\
\texttt{ 1234567} \\
\texttt{╔═══════╗} \\
\texttt{║W◌◌◌◌◌◌║ 1} \\
\texttt{║◌H◌◌◌◌◌║ 2} \\
\texttt{║T◌C◌◌┌─╢ 3} \\
\texttt{║I◌◌┌─┤◌║ 4} \\
\texttt{║◌◌◌│◌├─╢ 5} \\
\texttt{╟───┼─┘◌║ 6} \\
\texttt{║◌◌◌│◌◌◌║ 7} \\
\texttt{╚═══╧═══╝} \\
\texttt{```} \\
\texttt{Trenutno igrate krog 13 od največ 20.} \\
\texttt{Skupaj ste zbrali 11 od 12 kazni.} \\
\texttt{Drugi igralec je poslal to sporočilo:} \\
\\ 
\\ 
\texttt{"končano?"} \\
\\ 
\\ 
\texttt{Kakšen je vaš naslednji ukaz?} \\
            }
        }
    }
    & \\ \\

    \theutterance \stepcounter{utterance}  
    & \multicolumn{4}{p{0.6\linewidth}}{
        \cellcolor[rgb]{0.9,0.9,0.9}{
            \makecell[{{p{\linewidth}}}]{
                \texttt{\tiny{[P1$\rangle$GM]}}
                \texttt{RECI: končano!} \\
            }
        }
    }
    & & \\ \\

    \theutterance \stepcounter{utterance}  
    & & & \multicolumn{2}{p{0.3\linewidth}}{
        \cellcolor[rgb]{0.9,0.9,0.9}{
            \makecell[{{p{\linewidth}}}]{
                \texttt{\tiny{[GM$|$GM]}}
                \texttt{RECI: končano!} \\
            }
        }
    }
    & & \\ \\

    \theutterance \stepcounter{utterance}  
    & & & \multicolumn{2}{p{0.3\linewidth}}{
        \cellcolor[rgb]{0.9,0.9,0.9}{
            \makecell[{{p{\linewidth}}}]{
                \texttt{\tiny{[GM$|$GM]}}
                \texttt{true} \\
            }
        }
    }
    & & \\ \\

    \theutterance \stepcounter{utterance}  
    & & & \multicolumn{2}{p{0.3\linewidth}}{
        \cellcolor[rgb]{0.9,0.9,0.9}{
            \makecell[{{p{\linewidth}}}]{
                \texttt{\tiny{[GM$|$GM]}}
                \texttt{Initial states:} \\
\texttt{```} \\
\texttt{Player 1:} \\
\texttt{ 1234567} \\
\texttt{╔═══════╗} \\
\texttt{║◌◌◌◌◌◌◌║ 1} \\
\texttt{║◌◌WH◌◌◌║ 2} \\
\texttt{║◌C◌◌◌┌─╢ 3} \\
\texttt{║◌◌◌┌─┤◌║ 4} \\
\texttt{║◌◌I│◌├─╢ 5} \\
\texttt{╟───┼─┘◌║ 6} \\
\texttt{║◌◌T│◌◌◌║ 7} \\
\texttt{╚═══╧═══╝} \\
\\ 
\\ 
\texttt{Player 2:} \\
\texttt{ 1234567} \\
\texttt{╔═══════╗} \\
\texttt{║◌◌◌◌◌◌◌║ 1} \\
\texttt{║◌◌H◌◌◌◌║ 2} \\
\texttt{║◌WC◌◌┌─╢ 3} \\
\texttt{║◌◌◌┌─┤◌║ 4} \\
\texttt{║T◌◌│I├─╢ 5} \\
\texttt{╟───┼─┘◌║ 6} \\
\texttt{║◌◌◌│◌◌◌║ 7} \\
\texttt{╚═══╧═══╝} \\
\\ 
\texttt{```} \\
            }
        }
    }
    & & \\ \\

    \theutterance \stepcounter{utterance}  
    & & & \multicolumn{2}{p{0.3\linewidth}}{
        \cellcolor[rgb]{0.9,0.9,0.9}{
            \makecell[{{p{\linewidth}}}]{
                \texttt{\tiny{[GM$|$GM]}}
                \texttt{End states:} \\
\texttt{```} \\
\texttt{Player 1:} \\
\texttt{ 1234567} \\
\texttt{╔═══════╗} \\
\texttt{║W◌◌◌◌◌◌║ 1} \\
\texttt{║◌H◌◌◌◌◌║ 2} \\
\texttt{║T◌C◌◌┌─╢ 3} \\
\texttt{║I◌◌┌─┤◌║ 4} \\
\texttt{║◌◌◌│◌├─╢ 5} \\
\texttt{╟───┼─┘◌║ 6} \\
\texttt{║◌◌◌│◌◌◌║ 7} \\
\texttt{╚═══╧═══╝} \\
\\ 
\\ 
\texttt{Player 2:} \\
\texttt{ 1234567} \\
\texttt{╔═══════╗} \\
\texttt{║W◌◌◌◌◌◌║ 1} \\
\texttt{║◌H◌◌◌◌◌║ 2} \\
\texttt{║T◌C◌◌┌─╢ 3} \\
\texttt{║I◌◌┌─┤◌║ 4} \\
\texttt{║◌◌◌│◌├─╢ 5} \\
\texttt{╟───┼─┘◌║ 6} \\
\texttt{║◌◌◌│◌◌◌║ 7} \\
\texttt{╚═══╧═══╝} \\
\\ 
\texttt{```} \\
            }
        }
    }
    & & \\ \\

    \theutterance \stepcounter{utterance}  
    & & & \multicolumn{2}{p{0.3\linewidth}}{
        \cellcolor[rgb]{0.9,0.9,0.9}{
            \makecell[{{p{\linewidth}}}]{
                \texttt{\tiny{[GM$|$GM]}}
                \texttt{* success: True} \\
\texttt{* lose: False} \\
\texttt{* aborted: False} \\
\texttt{{-}{-}{-}{-}{-}{-}{-}} \\
\texttt{* Shifts: 6.00} \\
\texttt{* Max Shifts: 8.00} \\
\texttt{* Min Shifts: 4.00} \\
\texttt{* End Distance Sum: 0.00} \\
\texttt{* Init Distance Sum: 8.24} \\
\texttt{* Expected Distance Sum: 20.95} \\
\texttt{* Penalties: 11.00} \\
\texttt{* Max Penalties: 12.00} \\
\texttt{* Rounds: 13.00} \\
\texttt{* Max Rounds: 20.00} \\
\texttt{* Object Count: 5.00} \\
            }
        }
    }
    & & \\ \\

\end{supertabular}
}

\end{document}
